\documentclass[11pt]{article}
\usepackage[a4paper,margin=25mm]{geometry}
\usepackage{amsmath,amssymb,amsthm,booktabs,array,longtable}
\usepackage[T1]{fontenc}
\usepackage{lmodern}
\usepackage[colorlinks=true,urlcolor=blue,citecolor=blue,linkcolor=blue]{hyperref}
\newtheorem{theorem}{Theorem}
\newtheorem{corollary}{Corollary}
\newcommand{\ind}{\mathbf{1}}
\newcommand{\Z}{\mathbb{Z}}
\setlength{\emergencystretch}{2em}
\title{OEIS A240513, revision 11, FORMULA 1--3:\\
ladder proofs and extensions of Bala (2014)}
\author{Haobo Ma\thanks{ChronoAI Pte Ltd, Singapore.}\and
Wenlin Zhang\thanks{National University of Singapore; ChronoAI Pte Ltd, Singapore.}}
\date{11 November 2026}
\begin{document}
\maketitle
\begin{abstract}
We prove the three empirical formulas in OEIS A240513, revision 11,
FORMULA entries 1--3, and extend Bala's (2014) Chebyshev treatment of
the related norm sequence by making the ladder connection explicit
\cite{ladder,bala}. These formulas are Hardin's order-five recurrence,
Barker's generating function, and Mathar's relation to A105309.
The objects are binary two-column arrays in which no cell
has exactly two equal-valued orthogonal neighbors, with the first cell
fixed to zero. A transfer on sixteen ordered pairs of rows gives the
recurrence and generating function. A norm identity proves the relation
to A105309. The same construction proves the empirical width-three
and width-four recurrences in A240514--A240515 and A240519; their
minimal orders are 20 and 48. At every fixed width the generating
function is rational. We also give correctly indexed Chebyshev
identities and correct the sign and shift in A105309's
Chebyshev-transform formula.
\end{abstract}

\section{The array count}
The question in A240513 is whether its empirical recurrence, generating
function and quaternion relation hold for the stated binary arrays.
Theorem~\ref{thm:main} and Corollary~\ref{cor:norm} prove all three.
For $n\geq1$, let $a_n$ count maps
$A:\{1,\ldots,n\}\times\{1,2\}\longrightarrow\{0,1\}$ with $A(1,1)=0$
such that each cell has a number of equal-valued orthogonal neighbors
different from two. Neighbors must lie inside the rectangle: there is no
wrapping or diagonal adjacency. Rows and columns are labeled. In the
binary alphabet, introducing values in row-major order, as in the NAME
field of A240513, is equivalent to fixing the first cell to zero;
the value one need not occur. Set $a_0=0$ as a series convention, rather
than counting the empty rectangle. Write
\[
 Q(x)=1-x-2x^2-x^3+x^4,\qquad F(x)=\sum_{n\geq1}a_nx^n.
\]

\begin{theorem}\label{thm:main}
For every $n\geq6$,
\begin{equation}\label{eq:rec}
 a_n=2a_{n-1}+a_{n-2}-a_{n-3}-2a_{n-4}+a_{n-5}.
\end{equation}
Moreover, as a formal power series in $\Z[[x]]$,
\begin{equation}\label{eq:gf}
 F(x)=\frac{x(2-x)(1-x^2-2x^3)}{(1-x)Q(x)}.
\end{equation}
The initial values are $a_1,\ldots,a_6=2,3,6,10,21,42$.
\end{theorem}

These are FORMULA entries 1 and 2 of A240513, revision 11
(9 November 2018). The sequence was introduced by R.\ H.\ Hardin on
6 April 2014; the generating function is attributed to Colin Barker,
24 February 2018. The same recurrence appears in A240519, FORMULA,
$k=2$ \cite{rectangle}. All divisions below are valid formally because
the denominators have constant coefficient one.

\section{A complete transfer proof}
Let $\mathcal R=\{00,01,10,11\}$, with row coordinates indexed by $j=0,1$.
Put $e(u,v)=\ind_{u=v}$ and define the boundary and interior predicates
\begin{align*}
 B(r,s)&\iff e(r_j,r_{1-j})+e(r_j,s_j)\ne2\quad(j=0,1),\\
 I(r,s,t)&\iff e(s_j,s_{1-j})+e(s_j,r_j)+e(s_j,t_j)\ne2
                 \quad(j=0,1).
\end{align*}
Thus $B(r,s)$ checks a boundary row $r$ adjacent to $s$, whereas
$I(r,s,t)$ checks the middle row $s$. For a function
$w:\mathcal R^2\to\Z$, define
\[
 (Tw)(r,s)=\sum_{t\in\mathcal R}\ind_{I(r,s,t)}w(s,t),\qquad
 L(w)=\sum_{r,s\in\mathcal R}\ind_{r_0=0}\ind_{B(r,s)}w(r,s).
\]
Set $v_0(r,s)=\ind_{B(s,r)}$ and $v_{k+1}=Tv_k$.

For fixed initial rows $r,s$, the number of choices of $k$ further rows
that satisfy all interior conditions and the last boundary condition is
$v_k(r,s)$. Indeed, for $k=0$ the only condition is $B(s,r)$.
For $k+1$ further rows, choose the next row $t$, enforce $I(r,s,t)$,
and count the remaining $k$ rows by the induction hypothesis.
This gives exactly $Tv_k$. Enforcing the first boundary and first cell
therefore proves, for every $k\geq0$,
\begin{equation}\label{eq:bridge}
 a_{k+2}=L(v_k).
\end{equation}
Each array has exactly one row sequence, so no symmetry quotient or
multiplicity is introduced.

The following table lists all coordinates needed for the argument.
Each column after $v_0$ is obtained by the displayed definition of $T$.
\begin{center}
\small
\begin{tabular}{ccrrrrrr}
\toprule
$r$&$s$&$v_0$&$v_1$&$v_2$&$v_3$&$v_4$&$v_5$\\
\midrule
00&00&0&0&0&0&0&0\\
00&01&1&1&3&6&12&26\\
00&10&1&1&3&6&12&26\\
00&11&1&1&1&1&1&1\\
01&00&0&1&1&3&6&12\\
01&01&1&1&2&5&9&20\\
01&10&1&2&5&9&20&41\\
01&11&0&1&1&3&6&12\\
10&00&0&1&1&3&6&12\\
10&01&1&2&5&9&20&41\\
10&10&1&1&2&5&9&20\\
10&11&0&1&1&3&6&12\\
11&00&1&1&1&1&1&1\\
11&01&1&1&3&6&12&26\\
11&10&1&1&3&6&12&26\\
11&11&0&0&0&0&0&0\\
\bottomrule
\end{tabular}
\end{center}
Checking each of these sixteen rows gives the vector identity
\begin{equation}\label{eq:certificate}
 v_5=2v_4+v_3-v_2-2v_1+v_0.
\end{equation}
Apply $T^k$ and then $L$. Both maps are linear over $\Z$, so
\eqref{eq:bridge} yields \eqref{eq:rec} for $n=k+7$, hence $n\geq7$.
The two one-row arrays give $a_1=2$; the values
$L(v_0),\ldots,L(v_4)$ are $3,6,10,21,42$.
At $n=6$, the right side of \eqref{eq:rec} is
$2\cdot21+10-6-2\cdot3+2=42$, establishing the remaining case.

Finally,
\[
 (1-x)Q(x)=1-2x-x^2+x^3+2x^4-x^5.
\]
The coefficients of degrees zero through five in $(1-x)Q(x)F(x)$
are $0,2,-1,-2,-3,2$, using $a_0=0$ and the initial values.
All higher coefficients vanish by \eqref{eq:rec}. Thus
\[
 (1-x)Q(x)F(x)=2x-x^2-2x^3-3x^4+2x^5
             =x(2-x)(1-x^2-2x^3),
\]
which proves \eqref{eq:gf} and Theorem~\ref{thm:main}.

\section{The third empirical formula and adjacent entries}
In A105309 \cite{norm}, let $q_0=q_1=1$ and
$q_n=q_{n-1}+q_{n-2}J$, where
$J=(0,1/\sqrt3,1/\sqrt3,1/\sqrt3)$ is a unit imaginary quaternion.
Because $J^2=-1$, the real subalgebra spanned by $1,J$ is isomorphic
to $\mathbb C$, preserving norms. Consequently A105309 is
$c_n=|z_n|^2$, where
\[
 z_0=z_1=1,\qquad z_n=z_{n-1}+i z_{n-2}.
\]

\begin{corollary}\label{cor:norm}
For every $n\geq1$, $a_n=1+c_n$. In particular, for $n\geq5$,
\[
 a_n=a_{n-1}+2a_{n-2}+a_{n-3}-a_{n-4}-2.
\]
\end{corollary}
\begin{proof}
For $n\geq1$, write $w_n=z_n\overline{z_{n-1}}=u_n+i t_n$.
The defining recursion gives, for $n\geq2$,
\[
 w_n=c_{n-1}+i\overline{w_{n-1}},\quad
 u_n=c_{n-1}+t_{n-1},\quad t_n=u_{n-1},\quad
 c_n=c_{n-1}+c_{n-2}+2t_{n-1}.
\]
For $n\geq4$, these equalities imply
\begin{align*}
 c_n-c_{n-1}-c_{n-2}
 &=2u_{n-2}=2c_{n-3}+2t_{n-3}\\
 &=c_{n-2}+c_{n-3}-c_{n-4}.
\end{align*}
Together with $c_0,c_1,c_2,c_3=1,1,2,5$, coefficient comparison gives
\[
 C(x):=\sum_{n\geq0}c_nx^n=\frac{1-x^2}{Q(x)}.
\]
Direct substitution in \eqref{eq:gf} now yields
\[
 F(x)=\frac{x}{1-x}+C(x)-1.
\]
Comparing positive-degree coefficients proves the first assertion,
which is Mathar's FORMULA entry of 9 November 2018 in A240513.
Substituting $c_j=a_j-1$ in the order-four recurrence for $c_n$ proves
the second assertion when $n-4\geq1$.
\end{proof}

The generating function for $c_n$ is also recorded in A105309,
FORMULA (Mathar, 26 April 2024). A related earlier enumeration is
Mathar's non-bonding domino count \cite[p.~8, Eq.~(12)]{mathar}:
\[
 \widehat D_2(x,1)=\frac{1+x+x^2-x^3}{Q(x)}
                  =\frac{C(x)-1}{x}.
\]
This is the shifted sequence $c_{n+1}$, not the ladder sequence $a_n$.
The argument above establishes the ladder relation directly, without
requiring a domino-to-ladder bijection.

The rational function $xC(x)$ already occurs in Chen and Shapiro's
skinny-tree enumeration \cite[Theorem~16]{chen}, published in 2007.
Bala \cite[Eqs.~(5)--(6), (19)--(22)]{bala} treats the same shifted
sequence as a fourth-order divisibility sequence, with parameters
$(P_1,P_2,Q)=(1,-4,1)$. These are antecedents for the norm sequence;
the transfer proof identifies the binary-ladder count with it.

\begin{corollary}\label{cor:rectangle}
Let $T(n,k)$ be the rectangular-array count of A240519 with the same
local rule and first-cell convention. Then, for $n,k\geq1$,
\[
 T(n,k)=T(k,n),\qquad T(n,2)=T(2,n)=a_n,\qquad
 T(n,1)=T(1,n)=F_{n+1},
\]
where $F_1=F_2=1$ are Fibonacci numbers.
\end{corollary}
\begin{proof}
Transposition preserves the first cell and all orthogonal adjacencies.
The width-two assertion is the definition of $a_n$.
For width one, only interior cells can have two neighbors, and the
forbidden patterns are $000$ and $111$. Hence every constant run has
length one or two. Since the initial symbol is fixed and successive
runs alternate, admissible strings correspond bijectively to
compositions of $n$ into parts one and two. Their number is $F_{n+1}$,
by separating compositions according to their first part.
\end{proof}

The transfer description also extends to every fixed width $w\geq1$.
Let $\mathcal R_w=\{0,1\}^w$ and let $h_j(r)$ be the number of
in-row neighbors of position $j$ equal to $r_j$ (zero, one, or two).
Replace $B$ and $I$ by
\[
 B_w(r,s)\iff h_j(r)+e(r_j,s_j)\ne2\quad\text{for every }j,
\]
\[
 I_w(r,s,t)\iff h_j(s)+e(s_j,r_j)+e(s_j,t_j)\ne2
                    \quad\text{for every }j.
\]
Define $T_w,L_w,b_w$ by the same formulas as $T,L,v_0$, using
$\mathcal R_w,B_w,I_w$ and the first-coordinate constraint $r_0=0$.
The continuation induction in \eqref{eq:bridge} is unchanged, proving
$T(n,w)=L_w(T_w^{n-2}b_w)$ for $n\geq2$. Since $T_w$ is an integer
matrix on $4^w$ states, summing its powers proves
\[
 \sum_{n\geq1}T(n,w)x^n
 =F_{w+1}x+x^2 L_w\bigl((\mathrm{Id}-xT_w)^{-1}b_w\bigr).
\]
The inverse exists formally and has rational entries: its determinant
has constant coefficient one. Thus the count has a rational generating
function at every fixed width.

\begin{corollary}\label{cor:orders}
The columns $T(n,3)$ and $T(n,4)$ satisfy the empirical recurrences in
A240514 and A240515 \cite{wider}, respectively, for $n\geq21$ and $n\geq49$.
Their minimal eventual homogeneous recurrence orders over $\mathbb Q$
are 20 and 48. Thus both empirical orders in A240519, FORMULA $k=3,4$,
hold, with minimality as an additional conclusion.
\end{corollary}
\begin{proof}
Write $d_3=20$, $d_4=48$ and
$D_w(x)=\sum_{j=0}^{d_w}q_{w,j}x^j$.
In increasing degree order the coefficients are
\[
\begin{split}
(q_{3,j})={}&(1,-4,3,2,-2,-2,7,8,-15,-4,8,4,-15,-8,7,2,\\
           &\hspace{7mm}-2,-2,3,4,1),
\end{split}
\]
\[
\begin{split}
(q_{4,j})={}&(1,-5,5,-4,3,49,-10,-15,-45,-124,-131,-301,-18,\\
 &615,201,1508,823,413,254,-2467,-793,-4516,-1071,-2289,0,\\
 &2289,1071,4516,793,2467,-254,-413,-823,-1508,-201,-615,\\
 &18,301,131,124,45,15,10,-49,-3,4,-5,5,-1).
\end{split}
\]
These reproduce the coefficients of the two source recurrences after
moving all terms to the left. Put $v_{w,k}=T_w^k b_w$.
For $w=3,4$, the finite identities needed below are
\[
 \sum_{j=0}^{d_w}q_{w,j}v_{w,d_w-j}=0,
 \qquad \det H_w\not\equiv0\pmod {1000000007},
\]
where $H_w=(T(i+j+2,w))_{0\leq i,j<d_w}$.
The accompanying \texttt{verify.py} checks the vector identities using
precisely $B_w,I_w$ above and checks the determinants in exact modular
arithmetic. The state counts, determinant residues, coefficients and
initial values are in \texttt{width-certificates.json} and
\texttt{checks.json}. The identities do not depend on a fitted finite
prefix.

Applying $L_wT_w^k$ to the zero vector proves
$\sum_{j=0}^{d_w}q_{w,j}T(n-j,w)=0$ for $n\geq d_w+2$.
The remaining boundary case $n=d_w+1$ follows by substituting
$T(1,3)=3$ or $T(1,4)=5$ and the first $d_w$ transfer totals; the
same integer check gives zero. This proves both recurrences at their
first possible indices.

The nonzero modular determinants imply $\det H_w\ne0$ over
$\mathbb Q$. To see why this also proves eventual minimality,
let $g_k=T(k+2,w)$ and form the state
$(g_k,\ldots,g_{k+d_w-1})^{\mathsf t}$. Its shift is the companion
matrix of the recurrence, whose determinant is $\pm q_{w,d_w}=\pm1$.
Consequently the Hankel matrix beginning at $g_s$ is the product of
the $s$th power of this invertible companion matrix and $H_w$.
Every shifted Hankel matrix is therefore nonsingular. An eventual
recurrence of smaller order would make its sufficiently late $d_w$
consecutive columns dependent, a contradiction.
\end{proof}

\section{Further norm identities and their indexing}
Let $\alpha=(1+\sqrt{17})/4$, $\beta=(1-\sqrt{17})/4$,
$M=\left(\begin{smallmatrix}0&1\\1&1/2\end{smallmatrix}\right)$ and
$\gamma=(1+i)/\sqrt8$.
Use the usual recurrences $T_0=1,T_1=t$,
$T_{m+1}=2tT_m-T_{m-1}$ and $U_0=1,U_1=2t$,
$U_{m+1}=2tU_m-U_{m-1}$.

\begin{corollary}\label{cor:chebyshev}
For every $n\geq0$,
\[
 c_n=\frac{T_{n+1}(\alpha)-T_{n+1}(\beta)}{\alpha-\beta}
     =[T_{n+1}(M)]_{21}
     =U_n(\gamma)U_n(\overline\gamma).
\]
If $\mathcal C A(x)=\frac{1-x^2}{1+x^2}
A\bigl(\frac{x}{1+x^2}\bigr)$ denotes the Chebyshev transform,
then
\[
 \mathcal C\!\left(\frac{t}{1-t-4t^2}\right)(x)=xC(x).
\]
Thus $a_n$ is one plus each of the first three expressions for $n\geq1$.
\end{corollary}
\begin{proof}
The generating function of $T_{n+1}(t)$ is
$(t-x)/(1-2tx+x^2)$. Subtract at $\alpha,\beta$ and divide
by $\alpha-\beta$. Since $\alpha+\beta=1/2$ and
$\alpha\beta=-1$, the result is $(1-x^2)/Q(x)=C(x)$.
For every polynomial $p$, reduction modulo
$M^2-M/2-\mathrm{Id}=0$ writes $p(M)=A\mathrm{Id}+BM$.
Evaluation at the distinct roots $\alpha,\beta$ gives
$B=(p(\alpha)-p(\beta))/(\alpha-\beta)$, which is its lower-left
entry.
Set $\eta=(1-i)/\sqrt2$. The sequence
$\eta^nU_n(\gamma)$ has initial values 1,1 and recurrence
$z_n=z_{n-1}+iz_{n-2}$, because $2\eta\gamma=1$ and
$\eta^2=-i$. Its squared norm is $c_n$, proving the third expression.
Finally the substitution in $\mathcal C$ gives denominator
$(1+x^2)^2-x(1+x^2)-4x^2=Q(x)$ and numerator $x(1-x^2)$.
\end{proof}

The A105309 FORMULA block attributed to Bala (27 March 2014)
prints $T_n$, $T_n(M)$ and $U_{n-1}$ in place of
$T_{n+1}$, $T_{n+1}(M)$ and $U_n$. At $n=2$ all three
printed expressions equal one instead of $c_2=2$.
Its transform input has $+4t^2$ instead of $-4t^2$ and its output
requires the shift $xC(x)$, rather than $C(x)$.
These statements are refuted-and-corrected by
Corollary~\ref{cor:chebyshev}. The correctly indexed first-kind formula
already appears in Somos's Mathematica code in that entry
(20 December 2016). Bala's bivariate formula itself uses the shifted
sequence $W_{n+1}=c_n$ and is consistent with these corrections.
The corrected transform input is $t$ times the generating function
of A006131 \cite{auxiliary}.

\begin{corollary}\label{cor:auxiliary}
Let $r_n=\mathrm{A092886}(n)$, and let $p_n,s_n$ be A201837 and
A201838. Then
\[
 c_n=r_{n+1}-r_{n-1}\quad(n\geq1),\qquad
 c_n=p_n^2+s_n^2\quad(n\geq0).
\]
Extending $c_n$ by its order-four recurrence to all integers gives
$c_n=c_{-n-2}$. For $m,k\geq1$ one has the divisibility
$c_{m-1}\mid c_{mk-1}$; consequently
$a_{m-1}-1\mid a_{mk-1}-1$ when $m\geq2$.
\end{corollary}
\begin{proof}
The defining series for $r_n$ is $x/Q(x)$, so multiplication by
$1-x^2$ proves the first identity. The defining series for
$h_n=p_n+is_n$ is $(1-ix-ix^2)^{-1}$ \cite{auxiliary}.
Its recurrence gives $h_n=i^n\overline{z_n}$, with matching
initial values, and hence the second identity. Equivalently,
with $r_\pm=(i\pm\sqrt{4i-1})/2$,
\[
 h_n=\frac{r_+^{n+1}-r_-^{n+1}}{r_+-r_-},
\]
proving also Suteu's A105309 closed form (20 December 2016).

Put $u_0=0$, $u_m=h_{m-1}$ for $m\geq1$.
The recurrence $u_{m+1}=iu_m+iu_{m-1}$ extends to $\mathbb Z$.
Since $r_+r_-=-i$, its closed form gives
$u_{-m}=-(-i)^{-m}u_m$ and $|u_{-m}|^2=|u_m|^2$.
Thus $|u_{n+1}|^2$ has the stated reflection; it equals the
unique two-sided extension of $c_n$, since the order-four recurrence
has nonzero last coefficient and the norm calculation remains valid.
For divisibility, the identity
\[
 \begin{pmatrix}i&i\\1&0\end{pmatrix}^{m}
 =\begin{pmatrix}u_{m+1}&iu_m\\u_m&iu_{m-1}\end{pmatrix}
 \equiv iu_{m-1}\mathrm{Id}\pmod {u_m}
\]
shows, on raising to the $k$th power, that $u_m\mid u_{mk}$ in
$\mathbb Z[i]$. Taking norms proves the integer divisibility.
The last assertion uses $c_j=a_j-1$ at positive indices.
This rederives the shifted divisibility property in Bala's family
\cite[Eq.~(5)]{bala} and transfers it to the ladder count.
\end{proof}

\section{Coverage and verification}
\begin{longtable}{@{}p{0.43\linewidth}p{0.51\linewidth}@{}}
\toprule
Source item & Mathematical coverage\\
\midrule
A240513 rev.~11, FORMULA 1--2 & Proved by Theorem~\ref{thm:main}.\\
A240513 rev.~11, FORMULA 3 & Proved by Corollary~\ref{cor:norm}.\\
A105309, NAME and rational g.f. & Norm definition and generating function
proved in Corollary~\ref{cor:norm}; related shifted enumerations:
Chen--Shapiro (2007), Mathar (2024); prior proofs cited.\\
A240519, FORMULA $k=1,2$ & Proved by Corollary~\ref{cor:rectangle}
and Theorem~\ref{thm:main}.\\
A240519 rev.~6, FORMULA $k=3,4$; A240514--A240515, FORMULA &
Proved by Corollary~\ref{cor:orders}. Minimality: not posed; proved.\\
A240519, arbitrary fixed width & Rational generating function from the
$4^w$-state transfer; not posed as a formula, proved.\\
A105309, Bala FORMULA block (2014), first-/second-kind and matrix formulas &
Refuted-and-corrected in Corollary~\ref{cor:chebyshev}; corrected first-kind
formula also listed by Somos (2016).\\
A105309, Bala FORMULA block, transform & Refuted-and-corrected, with
input $t/(1-t-4t^2)$ and output $xC(x)$.\\
A105309, A092886 difference; Hanna's sum of squares (2011);
Suteu's closed form and Somos's reflection (2016) &
Proved by Corollary~\ref{cor:auxiliary}; matching auxiliary entries
A092886, A201837--A201838.\\
Bala (2014), divisibility family & Reproved; prior proof Bala (2014).
Ladder consequence in Corollary~\ref{cor:auxiliary}: not posed; proved.\\
Mathar (2024), unrelated larger-width domino and density statements &
Not covered; their objects differ from the binary-array model.\\
\bottomrule
\end{longtable}

\medskip
The pinned Lean file
\href{https://github.com/the-omega-institute/trureturing/blob/70badecc7fca07aa0b70e1202720702f7d1f9670/D5/S3/Combinatorics/TraceAvoidingArrays/HardinA240513Ladder.lean}{\texttt{HardinA240513Ladder.lean}}
\cite{lean} defines the literal array count
\texttt{a}. Its declaration \texttt{result} states both the recurrence
for $n\geq6$ and the
polynomial-multiplied identity \eqref{eq:gf} over $\Z[[x]]$;
\texttt{finite\_certificate} states \eqref{eq:certificate} on all sixteen
row-pair states. Corollaries~\ref{cor:norm}, \ref{cor:rectangle},
\ref{cor:chebyshev} and~\ref{cor:auxiliary}, and fixed-width rationality,
are written proofs. Corollary~\ref{cor:orders} combines the written
transfer and Hankel arguments with exhaustive finite integer and modular
checks. These additional results are outside the
scope of the pinned Lean declaration.
The accompanying \texttt{verify.py} checks the continuation-vector
identity, the wider-state annihilators and nonzero Hankel determinants.
It also compares literal arrays, recurrence coefficients, quaternion
norms and the related norm identities on the finite ranges specified
in \texttt{checks.json}; published OEIS tables provide optional further
comparisons. The exhaustive annihilator checks have the all-height
consequence proved in Corollary~\ref{cor:orders}; the prefix comparisons
are bounded corroboration. Full reproduction commands and finite ranges
are in the accompanying \texttt{README.md}.

\begin{thebibliography}{13}
\bibitem{ladder} R. H. Hardin, with formulas by Colin Barker and
R. J. Mathar, \emph{A240513}, The On-Line Encyclopedia of Integer Sequences,
revision 11, 9 November 2018. NAME, OFFSET, and FORMULA entries 1--3.
\url{https://oeis.org/A240513}.
\bibitem{norm} Gerald McGarvey and contributors, \emph{A105309},
The On-Line Encyclopedia of Integer Sequences, revision 70,
30 May 2026. NAME; FORMULA, generating function attributed to
R. J. Mathar, 26 April 2024. \url{https://oeis.org/A105309}.
\bibitem{rectangle} R. H. Hardin, \emph{A240519},
The On-Line Encyclopedia of Integer Sequences, revision 6,
23 July 2025. NAME and FORMULA, $k=1,2,3,4$.
\url{https://oeis.org/A240519}.
\bibitem{wider} R. H. Hardin, \emph{A240514} and \emph{A240515},
The On-Line Encyclopedia of Integer Sequences, NAME and empirical
FORMULA; introduced 6 April 2014.
\url{https://oeis.org/A240514}; \url{https://oeis.org/A240515}.
\bibitem{auxiliary} M. Somos, \emph{A092886}; P. D. Hanna,
\emph{A201837} and \emph{A201838}; N. J. A. Sloane and contributors,
\emph{A006131}, The On-Line Encyclopedia of Integer Sequences.
\url{https://oeis.org/A092886}; \url{https://oeis.org/A201837};
\url{https://oeis.org/A201838}; \url{https://oeis.org/A006131}.
\bibitem{chen} Ricky X. F. Chen and Louis W. Shapiro,
\emph{On Sequences $G_n$ Satisfying $G_n=(d+2)G_{n-1}-G_{n-2}$},
Journal of Integer Sequences \textbf{10} (2007), Article 07.8.1,
Theorem 16, p.~8. Published 2 August 2007.
\url{https://cs.uwaterloo.ca/journals/JIS/VOL10/Chen/chen509.pdf}.
\bibitem{bala} Peter Bala,
\emph{Linear Divisibility Sequences and Chebyshev Polynomials},
March 2014, Eqs.~(5)--(6), (19)--(22).
\url{https://oeis.org/A100047/a100047.pdf}.
\bibitem{mathar} Richard J. Mathar,
\emph{Bivariate Generating Functions Enumerating Non-Bonding Dominoes
on Rectangular Boards}, arXiv:2404.18806v1 (2024), p.~8, Eq.~(12).
\url{https://arxiv.org/abs/2404.18806v1}.
\bibitem{lean} \emph{HardinA240513Ladder.lean}, revision
\texttt{70badecc7fca07aa0b70e1202720702f7d1f9670}.
\href{https://github.com/the-omega-institute/trureturing/blob/70badecc7fca07aa0b70e1202720702f7d1f9670/D5/S3/Combinatorics/TraceAvoidingArrays/HardinA240513Ladder.lean}{Pinned Lean source: HardinA240513Ladder.lean}.
\end{thebibliography}

\noindent\small Copyright 2026 Haobo Ma and Wenlin Zhang.
Licensed under Creative Commons Attribution 4.0 International (CC BY 4.0).
\end{document}
